
This study addressed the question of which mechanism drives alignment-induced miscalibration in language models, providing the first pre-registered mechanistic discrimination among scale inflation (H1), decision-boundary restructuring (H2), and framing susceptibility (H3). The main findings, restricted to the Pythia 1.4B-6.9B family, are:

\begin{enumerate}
\item \textbf{H2 is the dominant mechanism:} All nine Spearman rho values fall below the H1 threshold of 0.90; eight of nine fall below 0.85. Alignment changes answer preferences rather than merely amplifying confidence magnitudes.
\end{enumerate}

\begin{enumerate}
\item \textbf{DPO >= PPO > SFT ordering (exploratory):} DPO produces larger calibration degradation than PPO in all three sizes, reversing the pre-registered prediction. The 2.8B case is inconclusive due to overlapping confidence intervals. This finding is based on public fallback checkpoints and requires matched-checkpoint replication.
\end{enumerate}

\begin{enumerate}
\item \textbf{H3 is not supported in softmax-ECE:} $\Delta ECE_TruthfulQA$ < $\Delta ECE_MMLU$ for all alignment types, indicating domain-general miscalibration in this evaluation framework.
\end{enumerate}

\begin{enumerate}
\item \textbf{Reusable framework:} The Spearman rho threshold test, argmax partition, and cross-benchmark diagnostic are applicable to any model family via lm-eval without modification.
\end{enumerate}

Three categories of follow-up work are motivated by these results: (a) matched-training DPO versus PPO experiments to verify the ordering observation; (b) testing whether ATS effectiveness correlates with the degree of H2 boundary shift (H-M4); and (c) applying the discrimination framework to larger and more widely deployed model families to determine whether H2 dominance persists at production scale.
